\section{Marked groups and intersections}
% BEGIN UF.Conventions
We use $x^g=g^{-1}xg$ and $[x,y]=x^{-1}y^{-1}xy$.
A rank-three string marking of a finite group $G$ is a generating triple
of distinct nonidentity involutions $r_0,r_1,r_2$ whose endpoints commute.
Write $G_I=\gen{r_i:i\in I}$ for $I\subseteq\{0,1,2\}$.
The marking is a \emph{string C-group} if
\[
 G_I\cap G_J=G_{I\cap J}\qquad(I,J\subseteq\{0,1,2\}).
\]
Its type is $\{k_1,k_2\}$, where $k_1=|r_0r_1|$ and
$k_2=|r_1r_2|$. We call it nondegenerate when $k_1,k_2\geq3$,
and tight when $|G|=2k_1k_2$. A marked isomorphism preserves the
three generators; duality reverses them. We retain the classical
correspondence between string C-groups and abstract regular polyhedra
when using geometric terminology.

% END UF.Conventions

% BEGIN UF.RankThree
\begin{lemma}[Rank-three intersection test]\label{UF.RankThree}
A nondegenerate rank-three string marking is a string C-group if and
only if
\begin{equation}\label{eq:rankthree}
 \gen{r_0,r_1}\cap\gen{r_1,r_2}=\gen{r_1}.
\end{equation}
In that case $2k_1k_2\leq |G|$.
\end{lemma}
\begin{proof}
Necessity is part of the intersection property. For sufficiency set
$H_{ij}=\gen{r_i,r_j}$. The endpoint group consists of
$1,r_0,r_2,r_0r_2$. If its intersection with $H_{01}$ contains $r_2$
or $r_0r_2$, then $r_2\in H_{01}\cap H_{12}=\gen{r_1}$,
contrary to distinctness. Hence
$H_{02}\cap H_{01}=\gen{r_0}$; dually
$H_{02}\cap H_{12}=\gen{r_2}$. These equalities and
\eqref{eq:rankthree} also exclude a generator from the group on the
other two. Singleton intersections follow, and pairs involving the
empty or full index set are immediate. This accounts for every pair
of subsets. Finally $|H_{01}|=2k_1$ and $|H_{12}|=2k_2$, so
\[
 |H_{01}H_{12}|=\frac{|H_{01}|\,|H_{12}|}{|H_{01}\cap H_{12}|}
 =2k_1k_2\leq |G|.
\]
\end{proof}
% END UF.RankThree

% BEGIN UF.PrimeIndex
\begin{lemma}[Prime-index intersection]\label{UF.PrimeIndex}
Let $C,H,J$ be subgroups of a group, with $C\leq H\cap J$ and
$[H:C]=p$ prime. If $H\not\leq J$, then $H\cap J=C$.
In particular, this holds if $\gen{H,J}=G$ and $J<G$.
\end{lemma}
\begin{proof}
The index tower gives
$p=[H:H\cap J][H\cap J:C]$. The first factor is not one,
so the second is one. In the final situation $H\leq J$ would give
$G=\gen{H,J}=J$, contrary to the hypothesis.
\end{proof}
For later products it is useful to record the elementary identities
\begin{equation}\label{eq:reflections}
 (sh)^2=1,\qquad (sh)(bh)=sb^{-1},
\end{equation}
whenever $h^2=1$, $hsh=s^{-1}$ and $hbh=b^{-1}$.
The first equality uses only the condition involving $s$; the second
uses only the condition involving $b$.
% END UF.PrimeIndex
